% ============================================================================
% 异常处理
% ============================================================================

\subsection{异常处理}

在 AArch64 架构中，异常是一个重要的机制。我们实现了完整的异常向量表和异常处理程序。

\subsubsection{异常向量表}

异常向量表要求 2048 字节对齐，包含 16 个入口：

\lstinputlisting[style=asm]{../QEMU/src/ASM/exceptions.s}

\subsubsection{异常处理 C 函数}

我们用 C 语言实现了异常解析和处理：

\lstinputlisting[style=cstyle]{../QEMU/src/exceptionhandler.c}

\subsubsection{\texttt{ESR\_EL2} 寄存器}

\texttt{ESR\_EL2} 寄存器存储异常原因，\texttt{EC} 字段（bit[31:26]）表示异常类别：

\begin{itemize}
    \item \texttt{0x0E}：非法执行状态
    \item \texttt{0x15}：SVC 指令（AArch64）
    \item \texttt{0x21}：指令中止（当前 EL）
    \item \texttt{0x25}：数据中止（当前 EL）
    \item 等等...
\end{itemize}

\subsubsection{工作流程}
\begin{enumerate}
    \item 异常发生，CPU 跳转到异常向量表
    \item 汇编保存所有通用寄存器
    \item 读取 \texttt{ESR\_EL2}、\texttt{ELR\_EL2}、\texttt{FAR\_EL2} 寄存器
    \item 调用 C 函数 \texttt{c\_exception\_handler} 解析和打印
    \item 如果是同步异常，返回 1 表示需要跳过触发指令
    \item 汇编恢复寄存器，执行 \texttt{eret} 返回
\end{enumerate}
